\documentclass{resume}

\begin{document}

\name{Ye Sun}
\contact{
Austin Robinson Building, Sidgwick Ave, Cambridge, UK, CB3 9DD \\
Email: \href{mailto:ys627@cam.ac.uk}{ys627@cam.ac.uk} \quad
Website: \href{https://ye-sun-econ.github.io/}{ye-sun-econ.github.io}
}

\begin{rSection}{Education}

\textbf{PhD in Economics}, University of Cambridge \hfill 2023–Present
\begin{rItemize}
    \item Supervisor: Vasco M. Carvalho \\ 
    Professor of Economics, University of Cambridge \\ 
    Email: \href{mailto:vmpmdc2@cam.ac.uk}{vmpmdc2@cam.ac.uk}
\end{rItemize}
\textbf{MA in Computational Social Sciences}, University of Chicago \hfill 2019–2021

\textbf{BA in Economics}, University of International Business and Economics \hfill 2015–2019

\end{rSection}

\begin{rSection}{Research Fields}
Macroeconomics, Production Networks, Firm Dynamics, International Trade
\end{rSection}

\begin{rSection}{Working Papers}


\textcolor{blue}{\textbf{Dynamic Endogenous Production Networks, Knowledge Diffusion, and Misallocation}}
\\
\small
\underline{Abstract:} I study how knowledge diffusion in supplier-customer relationships influences the structure and dynamics of production networks.
    I develop an optimal-stopping model that allows heterogeneous firms to endogenize the choice of staying with or switching away from the current supplier, 
    not based solely on cost minimization but on the potential for future gains through knowledge sharing. 
    The model predicts that higher rates of knowledge diffusion lead to a ``stickier'' production network. 
    Misallocation arises due to search frictions as firms are overly patient in forming long-term relationships, but knowledge diffusion mitigates its negative effects on aggregate productivity.
    Empirical analysis using data from Compustat's firm-level input-output data and USPTO patent data highlights the overlap between production and citation networks, showing that supplier-customer with citations tend to last longer and have bigger shares. Using an instrumental variable approach that leverages examiner leniency as an exogenous source of variation in patent citations, I find that supplier-customer relationships with knowledge diffusion experience significantly higher performance, indicating the existance of active learning. 

\normalsize
\href{https://hanbaeklee.github.io/Webpage/Lee_Sun_2026.pdf}{\textbf{Endogenous Plucking Through Networks: The Plucking Paradox}}{} (with Hanbaek Lee)
\\
\small
\underline{Abstract:} We study an \textit{endogenous plucking} mechanism in a parsimonious general equilibrium model where firms choose production linkages (network intensity). In expansions, high productivity encourages denser linkages, raising efficiency and output. Yet greater connectedness amplifies adverse TFP shocks, so downturns are disproportionately severe when recessions begin from a highly connected state. The model implies \textit{network-size} and \textit{duration dependence} of shock propagation. Due to a pecuniary externality, the decentralized equilibrium underinvests in network intensity. We characterize an optimal fiscal policy that decentralizes the constrained efficient allocation and delivers a \textit{Plucking Paradox}: the planner prefers higher output in normal times even though it increases exposure to rare, severe disasters.

\normalsize
\textcolor{blue}{\textbf{Dynamic Trade Imbalances in a Gravity Model under Incomplete Markets}}
\\
\small
\underline{Abstract:} TBA
\end{rSection}

\begin{rSection}{Teaching}

Teaching Assistant, University of Cambridge \hfill 2024–Present
\begin{rItemize}
    \item R201 Advanced Macroeconomics II \\
    PhD21 Computational Methods
\end{rItemize}
\end{rSection}

\begin{rSection}{Research Experience}
Research Assistant \hfill 2023–Present
\begin{rItemize}
    \item RA to Vasco M. Carvalho, Nir Jaimovich, and Sergio Rebelo
\end{rItemize}
Pre-doc, Sciences Po Paris \hfill 2021–2023
\begin{rItemize}
    \item RA to Johannes Boehm and Thomas Chaney
\end{rItemize}
Research Assistant, University of Chicago \hfill 2019–2021
\begin{rItemize}
    \item RA to Jonathan Dingel and Felix Tintelnot
\end{rItemize}


\end{rSection}

\begin{rSection}{Software}

\texttt{Coefplots.jl} — Publication-quality regression visualization in Julia

\texttt{GLFixedEffectModels.jl} — High-dimensional fixed effects estimation in Julia

\end{rSection}

\begin{rSection}{Programming Skills}
Julia, Python, STATA, \LaTeX, Git, KNITRO, Gurobi, GNU Make, Bash, MATLAB
\end{rSection}

\begin{rSection}{Presentations \& Workshops}
External conferences
\begin{rItemize}
    \item 2026: SED Conference, Athens University of Economics and Business
    \item 2026: Midwest Macro Meeting
\end{rItemize}
Internal seminars at University of Cambridge
\begin{rItemize}
    \item Networks Workshop; Trade and Spatial Economics Workshop; PhD Macro Workshop
\end{rItemize}

\end{rSection}
\begin{rSection}{Personal Information}

Nationality: Chinese

Languages: Chinese (native), English (fluent), French (intermediate)
\end{rSection}

\begin{rSection}{Funding \& Scholarships}
Faculty of Economics Trusts

The Janeway Institute Scholarship
\end{rSection}

\end{document}
